\usepackage{amsmath} % Para \dfrac
\usepackage{tensor} % Para \leftidx
\usepackage{tikz-3dplot}
\usepackage{cancel}
\usetikzlibrary{positioning}
\usetikzlibrary{patterns}
\usetikzlibrary{shapes.misc}
\usepackage{enumitem}
\setlist[enumerate,1]{label=\alph*)}
\usepackage[most]{tcolorbox} % Carga la librería con todas sus funciones
\newtcolorbox[auto counter, number within=section]{sol}[1][]{
  colback=black!5!white,     
  colframe=darkgray,   
  enhanced,                 
  breakable,                
  arc=2pt,
  % --- LÍNEAS PARA EL ÍNDICE ---
  title={Solución \thetcbcounter}, 
  list entry={Solución \thetcbcounter}, 
  add to list={toc}{subsubsection}, 
  % --- CORRECCIÓN DE ECUACIONES ---
  before upper={
    \setcounter{equation}{0}
    \renewcommand{\theequation}{\thetcbcounter.\arabic{equation}} % <--- USAR \thetcbcounter
  },
  #1                         
}                                    


\NewTColorBox[use counter from=sol]{cinematica}{ O{CadetBlue} m d"" !O{} }{
  enhanced,
  colframe=#1!75!black,
  colback=#1!5!white,
  fonttitle=\bfseries,
  title={#2},
  % Usamos una condición segura para el watermark
  code={\IfValueT{#3}{\tcbset{watermark text={#3}}}},
  #4
}

\newtcolorbox{form}[1][]{
  colback=red!5!white,      % Fondo rojo muy claro
  colframe=red!75!black,    % Borde rojo oscuro/quemado
  fonttitle=\bfseries,      
  coltitle=white,           % Título en blanco para mejor contraste con el rojo
  title={Fórmulas},         % Título por defecto
  enhanced,                 
  breakable,                                 % Bordes suaves
  before skip=10pt,         
  after skip=10pt,
  arc=2pt,
  #1                        
}
\newtcolorbox{marcos}[1][]{
  colback=CadetBlue!5!white,       % Fondo azul muy tenue
  colframe=CadetBlue!50!black,     % Borde azul marino oscuro
  fonttitle=\bfseries,      
  coltitle=white,             % Texto del título en blanco
  title={Marcos de referencia y sistemas de coordenadas}, % Título por defecto
  enhanced,                 
  breakable,                
  arc=2pt,                  
  before skip=10pt,         
  after skip=10pt,          
  #1                        
}

\newtcolorbox{cinematicaUnificada}[1]{
  enhanced,
  colback=white,
  colframe=gray!90!white,
  arc=2pt,
  fonttitle=\bfseries,
  coltitle=black,
  title={#1},
  top=4mm,
  separator sign none,
  breakable% Quita los dos puntos del título si fuera necesario
}

% Comando para los subtítulos internos con línea
\newcommand{\subtitulo}[1]{
    \tcbline % Línea divisoria de tcolorbox
    \nopagebreak
    \vspace{3pt}
    \begin{center}
        \textbf{\small --- #1 ---}
    \end{center}
    \vspace{3pt}
}


% Definición del estilo "Solución"

\usepackage{stackrel}
\usepackage{textcomp}
\newcommand{\transport} [3] {\leftidx{^\mathcal{#2}}{\dfrac{d}{dt}#1}=\leftidx{^\mathcal{#3}}{\dfrac{d}{dt}#1}+\velang{#2}{#3} \times #1}
\newcommand{\transporta} [3] {\leftidx{^\mathcal{#3}}{\dfrac{d}{dt}#1}+\velang{#2}{#3}\times #1}
\newcommand{\der}[2]{\leftidx{^\mathcal{#2}}{\dfrac{d}{dt}#1}}
\newcommand{\marc}[2]{\leftidx{^\mathcal{#2}}{#1}}

%velocidad y aceleracion angular
\newcommand{\velang}[2]{\stackrel{\mathcal{#1} \to \mathcal{#2}}{\omega}}
\newcommand{\accang}[2]{\stackrel{\mathcal{#1} \to \mathcal{#2}}{\alpha}}

\newcommand{\transportad} [3] {\underbrace{\der{}{#2}\left(#1\right)}_{1}+\underbrace{\velang{#2}{#3}\times #1}_{2}}
\newcommand{\transD} [5] {\der{}{#4}\underbrace{\left(#1\right)}_{#5}+\underbrace{#2}_{\velang{#3}{#4}}\times \underbrace{#1}_{#5}}

%%Comando de tikz

\newcommand{\Polaresxyz}[5]{
    % Argumentos:
    % #1: Ángulo 
    % #2: Letra sistema 1
    % #3: Letra sistema 2
    % #4: Color sistema 1
    % #5: Color sistema 2
    

    % 1. Círculo y Ejes principales (negro)
    \draw[thin] (0,0) circle (2.5);

    % 2. Vectores del Sistema Fijo (Azul)
    \draw[#4, ultra thick, -{Latex}] (0,0) -- (1.6,0) node[below right] {$\vec{#2}_x$};
    \draw[#4, ultra thick, -{Latex}] (0,0) -- (0,1.6) node[left] {$\vec{#2}_y$};

    % 3. Vectores del Sistema Móvil (Verde)
    \draw[#5, ultra thick, -{Latex}] (0,0) -- (35:1.6) node[above] {$\vec{#3}_x$};
    \draw[#5, ultra thick, -{Latex}] (0,0) -- (35+90:1.6) node[left] {$\vec{#3}_y$};

    % 4. Arco del ángulo y Punto B
    \draw[->] (0.9,0) arc (0:35:0.9) node[midway, right, xshift=2pt] {$#1$};
    \fill (35:2) circle (1.5pt);

    % 5. Etiqueta Central y círculo pequeño
    \draw (0,0) circle (0.15);
    \node[below left, xshift=-2pt, yshift=-2pt] at (0,0) {\textcolor{#4}{$\vec{#2}_z$} = \textcolor{#5}{$\vec{#3}_z$}};
    \fill (0,0) circle (1pt);
}
\newcommand{\Polaresrtz}[5]{
    % Argumentos:
    % #1: Ángulo 
    % #2: Letra sistema 1
    % #3: Letra sistema 2

    % 1. Círculo y Ejes principales (negro)
    \draw[thin] (0,0) circle (2.5);

    % 2. Vectores del Sistema Fijo (Azul)
    \draw[#4, ultra thick, -{Latex}] (0,0) -- (1.6,0) node[below right] {$\vec{#2}_x$};
    \draw[#4, ultra thick, -{Latex}] (0,0) -- (0,1.6) node[left] {$\vec{#2}_y$};

    % 3. Vectores del Sistema Móvil (Verde)
    \draw[#5, ultra thick, -{Latex}] (0,0) -- (35:1.6) node[above] {$\vec{#3}_r$};
    \draw[#5, ultra thick, -{Latex}] (0,0) -- (35+90:1.6) node[left] {$\vec{#3}_\theta$};

    % 4. Arco del ángulo y Punto B
    \draw[->] (0.9,0) arc (0:35:0.9) node[midway, right, xshift=2pt] {$#1$};
    \fill (35:2) circle (1.5pt);

    % 5. Etiqueta Central y círculo pequeño
    \draw (0,0) circle (0.15);
    \node[below left, xshift=-2pt, yshift=-2pt] at (0,0) {$\vec{#2}_z = \vec{#3}_z$};
    \fill (0,0) circle (1pt);
}
\newcommand{\PolaresGen}[5]{
    % Argumentos:
    % #1: Ángulo 
    % #2: Letra sistema 1
    % #3: Letra sistema 2

    % 1. Círculo y Ejes principales (negro)
    \draw[thin] (0,0) circle (2);

    % 2. Vectores del Sistema Fijo (Azul)
    \draw[#4, ultra thick, -{Latex}] (0,0) -- (1.6,0) node[below right] {$\vec{#2}_x$};
    \draw[#4, ultra thick, -{Latex}] (0,0) -- (0,1.6) node[left] {$\vec{#2}_y$};

    % 3. Vectores del Sistema Móvil (Verde)
    \draw[#5, ultra thick, -{Latex}] (0,0) -- (35:1.6) node[above] {$\vec{#3}_r$};
    \draw[#5, ultra thick, -{Latex}] (0,0) -- (35+90:1.6) node[left] {$\vec{#3}_\theta$};

    % 4. Arco del ángulo y Punto B
    \draw[->] (0.9,0) arc (0:35:0.9) node[midway, right, xshift=2pt] {$#1$};
    \fill (35:1.8) circle (1.5pt);

    % 5. Etiqueta Central y círculo pequeño
    \draw (0,0) circle (0.15);
    \node[below left, xshift=-2pt, yshift=-2pt] at (0,0) {$\vec{#2}_z = \vec{#3}_z$};
    \fill (0,0) circle (1pt);
}
